\section{Related Work}
\label{sec:related}

\paragraph{Masking augmentations.}
Cutout~\citep{devries2017cutout} removes a fixed-size square at a random
location and was among the first to show that occluding contiguous regions acts
as an effective regulariser.  Random Erasing~\citep{zhong2020randomerasing}
generalises this to rectangles of random aspect ratio and area, optionally
filled with noise or the dataset mean.  Both choose the masked region without
regard to image content.  ICD can be viewed as a content-aware Cutout: the
shape of the occlusion follows the model's saliency rather than a random box,
and the special case of solid zero-fill on the high-saliency tiles recovers a
saliency-targeted form of classical Cutout.  CutMix~\citep{yun2019cutmix}
replaces a region with a patch from a second image and mixes the labels
accordingly; unlike CutMix, ICD and AICD are single-image operations and leave
the label unchanged.

\paragraph{Saliency-aware augmentation.}
The closest line of work conditions augmentation on a saliency or attention
signal.  KeepAugment~\citep{gong2021keepaugment} computes a saliency map and
\emph{preserves} the most important region, restricting augmentation to the
remainder, with the goal of not destroying discriminative information.  ICD
adopts the opposite stance: it deliberately occludes the high-saliency region
to discourage over-reliance on it, while AICD occupies the same regime as
KeepAugment by perturbing only the low-saliency background.  The two methods
thus bracket the design space that KeepAugment occupies on one side.  Other
saliency-driven approaches operate through mixing rather than dropout:
SaliencyMix~\citep{uddin2021saliencymix} uses a saliency map to choose the
source patch in a CutMix-style operation, and SnapMix~\citep{huang2021snapmix}
uses class activation maps to assign mixed-label proportions in fine-grained
recognition.  These methods require a second image and modify the target;
ICD and AICD do neither.

\paragraph{Class activation maps.}
ICD and AICD rely on a saliency map to score image regions.  We use
GradCAM~\citep{selvaraju2017gradcam}, which weights the activations of a target
convolutional layer by the gradient of the predicted class score and produces a
coarse localisation map.  Numerous refinements exist, including
GradCAM\texttt{++}~\citep{chattopadhay2018gradcampp} for improved localisation
of multiple instances and gradient-free variants such as
EigenCAM~\citep{muhammad2020eigencam}, and any of these may in principle drive
the masking.  Our reference implementation builds on the
\texttt{pytorch-grad-cam} library~\citep{gildenblat2021gradcam}, which unifies
these methods behind a common interface.

\paragraph{Automated augmentation policies.}
A parallel body of work searches for augmentation policies rather than
conditioning on the model state.  AutoAugment~\citep{cubuk2019autoaugment}
learns a policy by reinforcement learning;
RandAugment~\citep{cubuk2020randaugment} reduces the search to two scalar
hyperparameters; and TrivialAugment~\citep{muller2021trivialaugment} dispenses
with search entirely, sampling a single transformation per image uniformly.
These methods are complementary to ICD and AICD: they define \emph{which}
photometric or geometric transformation to apply, whereas ICD and AICD define
\emph{where} to apply a content-aware occlusion.  The two can be composed.
